\documentclass[AMS,STIX2COL]{WileyNJD-v2}

\articletype{Research Article}%

\received{26 March 2026}
\revised{}
\accepted{}

\raggedbottom

\begin{document}

\title{Sparsity, Steps, and Silicon: Understanding Efficient Generative AI Inference Beyond Quantization}

\author[1]{Vrathik Shenoy K}

\author[2]{Dr. Amrithkala M Shetty}

\authormark{VRATHIK SHENOY K}

\address[1]{\orgdiv{Department of Computer Applications}, \orgname{Nitte Institute of Professional Education}, \orgaddress{\state{Mangalore, Karnataka}, \country{India}}}

\corres{*Vrathik Shenoy K, Department of Computer Applications, Nitte Institute of Professional Education, Mangalore. \email{shenoyvrathik@gmail.com}}

\abstract[Abstract]{Energy use during generative AI inference has turned into a real sustainability problem now that large language models (LLMs) and diffusion image generators run at scale. Quantization, which lowers numerical precision, is usually the first lever engineers reach for. Yet results from recent work show that the benefits are uneven and sometimes negative. In this paper we benchmark energy-efficient inference across two model families: LLMs (LLaMA~3, GPT-OSS) and diffusion models (Stable Diffusion, Flux). Full-precision baselines are compared against INT8, INT4, and GGUF variants on NVIDIA A100, H100, and RTX~4090 GPUs, with energy recorded alongside output quality. Three results stand out from the data. Sparse Mixture-of-Experts (MoE) models use roughly 4--5$\times$ less energy per token than dense models of similar quality, and the margin grows with newer MoE designs: GLM-5 activates only 5.9\% of its parameters per token, MiniMax~M2.5 just 4.3\%. Plain INT8 quantization of U-Net diffusion models, by contrast, can \textit{raise} energy use by as much as 64\%, while cutting inference steps via distillation saves far more. Hardware generation is the third lever: NVIDIA Blackwell now offers 25--50$\times$ the throughput per megawatt of Hopper, which is well beyond what quantization alone can achieve. Read alongside the 2024--2026 wave of new models and accelerators, these results suggest that sparse architectures, fewer inference steps, and better silicon together provide a more dependable route to sustainable generative AI than chasing lower numerical precision on its own.}

\keywords{Generative AI, Energy Efficiency, Mixture-of-Experts, Inference Optimization, Quantization, Diffusion Models, Green AI, Hardware Accelerators}

\jnlcitation{\cname{%
\author{Vrathik Shenoy K.}} (\cyear{2026}),
\ctitle{Sparsity, Steps, and Silicon: Understanding Efficient Generative AI Inference Beyond Quantization}, \cjournal{Expert Systems}, \cvol{2026;00:1--16}.}

\maketitle

%% ========== INTRODUCTION ==========
\section{Introduction}\label{sec1}

Generative artificial intelligence has moved from a research curiosity to an infrastructure-scale technology in a very short span of time. Large language models now serve hundreds of millions of queries per day, and text-to-image systems produce billions of images per month across commercial APIs and open-source pipelines. One consequence of this growth is easy to miss in the excitement around new capabilities: energy use. The International Energy Agency projects that AI-related data centre electricity demand in the United States alone could reach 426~TWh per year by 2030, almost three times the 2024 figure \cite{iea2024electricity}. Per-query numbers look small on their own. A single LLM completion uses around 0.3--2~Wh, and a single diffusion-model image costs roughly 0.15--0.85~Wh. Multiplied across billions of daily interactions, the totals become very large. Weissenbacher et al.\ measured a single ChatGPT conversation at about 0.05~kWh, around ten times the energy of a traditional search query \cite{weissenbacher2025energy}. That gives a sense of the scale.

The default response from the systems community has been quantization. The idea is simple: store weights and activations at lower precision (INT8, INT4, or mixed formats) so that each operation uses less silicon area and, in principle, less power. Quantization has a long history in efficient neural-network deployment \cite{dettmers2022gptint8, frantar2023gptq, lin2024awq}, and it is still useful. The problem is that recent evidence suggests quantization alone gives an incomplete picture of inference efficiency, and sometimes a misleading one. For diffusion models, naive INT8 quantization can \textit{increase} energy use by as much as 64\% due to runtime kernel overhead and precision-conversion costs \cite{huang2024tfmqdm, bertazzini2025hidden}. At the same time, architectural ideas like Mixture-of-Experts (MoE) routing produce much larger and more consistent savings by activating only part of a model's parameters per forward pass \cite{shazeer2017moe, fedus2022switch, lopez2025moe}. For diffusion models specifically, step-reduction methods (distillation in particular) cut energy roughly in proportion to the number of sampling iterations removed \cite{salimans2022progressive, meng2023distillation}. Quantization simply cannot deliver that kind of saving.

So the issue is really one of emphasis. The existing literature offers a lot of partial views: energy numbers for specific model families \cite{samsi2023words, ferreira2025prompts}, quantization ablations on a particular GPU \cite{liu2024survey}, or vendor data sheets that report peak FLOPS without saying anything about real workloads. What is still missing is a single benchmarking study that brings architecture (dense versus sparse), precision (FP16, INT8, INT4, GGUF), step count for diffusion models, hardware generation (A100 up through Blackwell), and the recent 2024--2026 wave of MoE LLMs, distilled diffusion models, and purpose-built accelerators into one comparative frame. Without something like that, practitioners do not have the evidence to make good deployment decisions. They risk spending weeks on quantization tweaks that save very little, and overlooking architectural choices that could cut energy by an order of magnitude.

Earlier work has filled in pieces of this picture. Strubell et al.\ started the ``Green AI'' conversation by putting a number on the training cost of large NLP models, although their focus stayed on training and not inference \cite{strubell2019energy}. Luccioni et al.\ moved the conversation toward deployment by reporting per-task energy for BLOOM and similar models, but they did not vary architecture or quantization in a systematic way \cite{luccioni2023power, luccioni2024power}. Patterson et al.\ argued that hardware efficiency gains would eventually slow down training emissions, but their analysis came before the MoE and distillation breakthroughs of 2024--2025 \cite{patterson2022carbon}. Liu et al.\ later produced a thorough survey of LLM quantization that recorded perplexity inflation of 50--100\% under naive W4A4 settings \cite{liu2024survey}. Huang et al.\ showed that temporal-feature-aware quantization (TFMQ-DM) is needed to keep diffusion-model output usable at low bit widths \cite{huang2024tfmqdm}. Among diffusion-energy studies, the most thorough is Bertazzini et al.\ with 17 models, although that work did not cover LLMs or current hardware \cite{bertazzini2025hidden}. No study so far has measured energy across both LLM and diffusion families under controlled conditions and then placed those results against the latest accelerators. Our paper sets out to do exactly that.

We frame the investigation around \textit{compute--energy proportionality}: the rough idea that, with everything else held constant, reducing active compute should reduce energy. This proportionality fails when quantisation kernels are poorly tuned, or when the workload becomes memory-bound. Real efficiency then drifts away from what theory would predict. Identifying when and why this happens is part of what we set out to do.

Five specific objectives guide the work:
\begin{enumerate}
\item Measure how INT8, INT4, and GGUF quantisation affect inference energy and output quality for LLaMA~3, GPT-OSS, Stable Diffusion, and Flux.
\item Quantify the energy advantage of sparse MoE architectures over dense models of similar capability.
\item Characterise the effect of hardware-level settings (DVFS, batch size) on inference energy efficiency.
\item Look at how output size (token count or image resolution) relates to total energy use.
\item Pin down the conditions under which quantisation actually \textit{raises} energy, and explain why.
\end{enumerate}

The contribution is both academic and practical. On the academic side, the paper offers a single comparative framework that future studies can build on as new models and hardware appear. On the practical side, the results help system architects and ML engineers focus on the optimisation levers that actually matter, which saves both energy and the engineering time that would otherwise go into ineffective quantisation work. At one million daily queries, switching from a dense model to an MoE equivalent saves an estimated 70,000--160,000~Wh per day. That works out to roughly 10--23 metric tonnes of CO$_2$ per year at the US average grid carbon intensity.

The remainder of this paper is organised using the Create a Research Space (CARS) method \cite{schwartz2020green}: Section~\ref{sec2} establishes the territory through a critical literature review; Section~\ref{sec3} details the benchmarking methodology; Section~\ref{sec4} occupies the niche with experimental results and an updated 2024--2026 landscape; Section~\ref{sec5} provides a critical discussion; and Section~\ref{sec6} concludes with contributions and future directions.


%% ========== LITERATURE REVIEW ==========
\section{Literature Review}\label{sec2}

The energy efficiency of generative AI inference sits at the intersection of three research threads: model architecture design, numerical-precision reduction, and hardware-system optimisation. The review below works through each thread, points out where the literature agrees and where it does not, and isolates the gap our study addresses. The discussion stays tied to the five objectives listed in Section~\ref{sec1}.

\subsection{Energy Profiling of LLM Inference}

Interest in inference energy accelerated once generative models transitioned from research artifacts to production services. Strubell et al.\ were among the first to sound the alarm, estimating that training a large Transformer emitted as much carbon as five cars over their lifetimes; their work, however, focused on training costs and predated the current inference-dominated era \cite{strubell2019energy}. Luccioni et al.\ shifted attention to deployment by profiling BLOOM-176B inference, reporting 0.004~kWh per query and arguing that per-query energy must become a standard reporting metric \cite{luccioni2023power}. Their subsequent study broadened this to multimodal models, finding that image-generation tasks consumed 10--100$\times$ more energy per request than text generation, though the comparison mixed different model families without controlling for capability \cite{luccioni2024power}.

Samsi et al.\ produced the first systematic GPU-level power measurements for LLM inference at scale, setting baseline figures on A100 hardware that later studies have used as a reference \cite{samsi2023words}. Ferreira et al.\ added measurements for further models, including Qwen2.5 variants, and the energy hierarchy from earlier work held up \cite{ferreira2025prompts}. Community benchmarks have since become more detailed. The ML.ENERGY Leaderboard publishes joules-per-token figures for dozens of LLMs, and MoE models consistently sit on the Pareto frontier of quality against energy \cite{mlenergy2025}. Niu et al.'s TokenPowerBench reports per-token power figures along similar lines: on an A100, GPT-OSS-20B (MoE with 3.6B active parameters) draws about 2.0~Wh per 100 tokens, while the dense LLaMA~3-70B uses 9.7~Wh \cite{niu2025tokenpowerbench}. The 4.8$\times$ gap is larger than raw FLOP counts predict, and reflects the memory-bandwidth savings that come with sparse routing. Marrou et al.\ added another variable to the picture, showing that the choice between greedy decoding, beam search, and sampling has a measurable effect on GPU energy use \cite{marrou2025decoding}.

A limitation common to these benchmarking efforts is their snapshot nature: results are hardware-specific and do not account for the rapid pace at which new MoE models are pushing activation ratios even lower.

\subsection{Quantisation and Its Limits}

Quantisation has been the workhorse of efficient deployment since the early days of edge ML. Dettmers et al.\ demonstrated that 8-bit matrix multiplication with mixed-precision decomposition could run 175B-parameter models with virtually no accuracy loss \cite{dettmers2022gptint8}. Frantar et al.\ advanced this with GPTQ, enabling accurate one-shot INT4 quantisation through an approximate second-order method \cite{frantar2023gptq}. Lin et al.\ introduced activation-aware weight quantisation (AWQ), showing that protecting a small fraction of salient weights during quantisation preserves quality better than uniform rounding \cite{lin2024awq}.

These successes have not always carried over into energy savings. Liu et al.\ reviewed more than 40 quantisation methods and concluded that \textit{quality-neutral} quantisation below 8 bits is still an open problem for very large models. W4A4 schemes can inflate perplexity by 50--100\% unless careful calibration, layer-wise sensitivity analysis, and sometimes quantisation-aware fine-tuning are added \cite{liu2024survey}. The other point worth noting is that energy savings from quantisation usually lag throughput gains. Quantised kernels add their own overhead: dequantisation shims, padding for alignment, and fallbacks to higher precision for sensitive operations such as softmax and layer normalisation. In our own runs, INT4 quantisation of LLaMA~3-7B \textit{raised} energy per token by 2\%, which gives a sense of how far realised efficiency can drift from the theoretical estimate.

The picture is worse for diffusion models. Huang et al.\ pointed out that naive post-training quantisation breaks the temporal feature consistency that iterative denoising depends on, producing visible artifacts even at INT8 and making INT4 unworkable without targeted handling \cite{huang2024tfmqdm}. Their TFMQ-DM framework brought quality back at W4A8 and saved about 13\% energy on Stable Diffusion~v1.5, which is useful but modest. Tseng et al.\ proposed QuEST, which pairs selective fine-tuning with low-bit quantisation to push the quality-energy frontier a bit further \cite{tseng2024quest}. Li et al.\ pushed to the other extreme with 1.58-bit quantisation for Flux models, mostly to map out the boundary of how far precision reduction can go \cite{li2024158bitflux}. None of these approaches matches the savings that come from simply running fewer denoising steps, as we show in Section~\ref{sec4}.

\subsection{Step Reduction and Distillation}

A parallel line of work has attacked diffusion-model efficiency from the sampling side. Salimans and Ho introduced progressive distillation, halving the required sampling steps with each distillation round \cite{salimans2022progressive}. Meng et al.\ extended this to guided diffusion, achieving 4-step generation with minimal quality loss \cite{meng2023distillation}. The energy implications are substantial: since each step involves a full forward pass through the denoising network, reducing steps from 50 to 4 cuts compute by roughly 12.5$\times$, with energy savings that are nearly proportional because the workload is compute-bound at typical resolutions.

Recent models have built these lessons into their design. Flux Schnell drops inference down to 4 steps through changes in architecture and training \cite{bfl2024flux}, and Z-Image-Turbo reaches 8-step generation by distilling a 6B-parameter single-stream DiT \cite{tongyi2025zimage}. Wu et al.'s Qwen-Image-2.0 packs generation and editing into a single 7B model. That is a 65\% parameter reduction from its 20B predecessor, and the model still tops the AI Arena ELO leaderboard. Its Lightning distilled variant reaches 4-step inference with a 42.55$\times$ speedup \cite{wu2025qwenimage}. These step-reduced models now sit at the top of efficiency leaderboards, which supports the broader point that \textit{step reduction beats quantisation as an energy-saving strategy for diffusion models}. That is the result we test in Objective~5.

\subsection{Hardware Acceleration and Generational Efficiency}

Hardware is the third efficiency lever, and it is usually discussed in isolation from algorithmic optimisation. NVIDIA's A100 and H100 GPUs were the baseline hardware for most inference benchmarks in 2023--2024 \cite{nvidia2023a100, nvidia2024h100}. DVFS tuning on these platforms has been reported to lift tokens-per-watt by 30--50\% with no observable quality cost \cite{mei2017dvfs, maliakel2025dvfs}. That is essentially a free win, but academic studies often skip past it. Yang et al.\ argued for total system power measurement (CPU, memory, and networking included) rather than GPU-only numbers, since GPU power is only 65--80\% of total system draw during inference \cite{yang2025system}.

The 2024--2026 hardware landscape introduces qualitative shifts that alter the entire efficiency calculus. NVIDIA's Blackwell B200 claims 25$\times$ per-token efficiency over Hopper, driven by FP4 precision support and a redesigned Transformer Engine \cite{nvidia2024b200}. The GB300 NVL72 pushes this to 50$\times$ throughput per megawatt, achieving an estimated 140$\times$ improvement in joules per token over A100 FP16 \cite{nvidia2025gb300}. Google's Ironwood TPU~v7 reports 29.3~TFLOPS/W, 2$\times$ better per-watt than TPU~v6e \cite{google2025ironwood}, and Cerebras's wafer-scale CS-3 eliminates off-chip memory bottlenecks entirely, achieving 2,100 tokens per second on LLaMA~70B at a system power of 23~kW---approximately 91~tok/s/kW \cite{cerebras2024cs3}. AMD's MI325X competes with 256~GB HBM3E and 1.31 TFLOPS/W at FP16 \cite{amd2024mi325x}. MLPerf Inference v5.0 results provide the first standardised cross-vendor comparison including these new accelerators \cite{mlperf2025inference}.

These advances mean that hardware generation selection now has a larger impact on inference energy than most algorithmic optimisations. Yet no prior study has situated model-level efficiency findings against this rapidly evolving hardware context. Our work addresses this gap.

\subsection{Summary and Gap Identification}

In summary, existing literature provides strong evidence that (a)~MoE sparsity outperforms quantisation for LLM energy savings \cite{lopez2025moe}, (b)~step reduction dominates quantisation for diffusion models \cite{bertazzini2025hidden}, and (c)~hardware generational leaps can dwarf both \cite{nvidia2025gb300}. What is absent is a unified study that benchmarks all three levers together, incorporates the latest 2024--2026 models and accelerators, and provides practitioners with an integrated decision framework. This paper fills that gap.


%% ========== METHODS ==========
\section{Methods}\label{sec3}

This study adopts a benchmark-based comparative analysis design, combining controlled experiments on standardised hardware with a structured synthesis of published performance data for models and accelerators released between 2024 and 2026. The choice of a benchmarking methodology, rather than a purely survey-based approach, is motivated by the need to generate directly comparable energy measurements under identical conditions---something that is not possible when aggregating results from disparate studies that vary in hardware, software stack, ambient temperature, and measurement granularity \cite{yang2025system}.

\subsection{Hardware Platforms and Configuration}

All primary experiments were conducted on three NVIDIA GPU platforms representing the spectrum from consumer to data-centre hardware: the RTX~4090 (consumer, 450~W TDP, 24~GB GDDR6X), the A100~80GB SXM (data-centre, 400~W TDP, 80~GB HBM2e, AMD EPYC~7763 host), and the H100~SXM (data-centre, 700~W TDP, 80~GB HBM3). Each system ran Ubuntu~22.04 with CUDA~12.2, PyTorch~2.3, and the latest vendor-optimised inference runtimes: vLLM for LLMs and diffusers with xFormers for image-generation models. To minimise confounding variables, all systems were isolated from other workloads during benchmarking runs, with idle power consumption measured and subtracted from active-inference power readings.

Dynamic Voltage and Frequency Scaling (DVFS) experiments varied the GPU core clock from the manufacturer's base frequency to the maximum boost frequency in 100~MHz increments, following the protocol of Maliakel et al.\ \cite{maliakel2025dvfs}. Batch sizes of 1, 4, 8, and 16 were tested for LLMs, while diffusion models were run exclusively at batch size~1, reflecting their typical deployment pattern.

\subsection{Models Under Evaluation}

For large language models, we evaluated LLaMA~3 in 7B and 70B dense variants \cite{grattafiori2024llama3} and GPT-OSS in 20B and 120B sparse MoE variants. Each model was tested at FP16 baseline and in INT8, INT4, and GGUF (Q8, Q4) quantised configurations. The GGUF variants were run through llama.cpp's optimised C++ inference engine, while INT8 and INT4 variants used the bitsandbytes library integrated with vLLM. Text-generation benchmarks consisted of 100 prompts drawn from the Open LLM Leaderboard evaluation suite, each generating 256 output tokens, with results averaged over three runs.

For diffusion models, we benchmarked Stable Diffusion v1.5, SD-XL, and SD~3.5~Medium (all U-Net-based), alongside Flux Dev and Flux Schnell (Transformer-based) \cite{rombach2022ldm, bfl2024flux}. Each model generated 100 images at 512$\times$512 resolution using its default sampling schedule. FP16 and INT8 variants were compared, with additional tests of TFMQ-DM (W4A8) quantisation on SD~v1.5 \cite{huang2024tfmqdm}.

\subsection{Energy Measurement Protocol}

GPU power was sampled at 100~ms intervals using \texttt{nvidia-smi dmon} throughout each inference run. Total energy per inference was computed as the integral of instantaneous power over wall-clock inference time, minus pre-measured idle power. All measurements were repeated at least three times, and we report medians with interquartile ranges. On the A100 and H100 platforms, we additionally used CodeCarbon at 1~Hz over 100 active runs (with 10 warm-up runs and idle power subtracted), following the methodology of Niu et al.\ \cite{niu2025tokenpowerbench}. Total system power (CPU + GPU + memory + networking) was recorded for a subset of runs using IPMI out-of-band readings, confirming that GPU power constitutes 65--80\% of total system power during inference-dominated workloads, consistent with Yang et al.'s findings \cite{yang2025system}.

\subsection{Quality Metrics}

For LLMs, output quality was assessed via perplexity on a held-out subset of WikiText-103 and task accuracy on MMLU and HellaSwag. For diffusion models, we used Fr\'{e}chet Inception Distance (FID) computed against a 5,000-image reference set from COCO~2017 validation, and CLIP score for text-image alignment, following the protocol of Bertazzini et al.\ \cite{bertazzini2025hidden}. A quantisation scheme was considered quality-neutral if perplexity increased by less than 1\% (LLMs) or FID increased by less than 5 points (diffusion models).

\subsection{Extended Analysis: 2024--2026 Models and Hardware}

Beyond our controlled experiments, we conducted a structured analysis of published benchmarks for the latest models and hardware released between 2024 and 2026. For LLMs, this includes GLM-5 (744B total, 40B active) \cite{zhipu2026glm5}, MiniMax~M2.5 (230B/10B) \cite{minimax2026m25}, LFM2-24B-A2B (24B/2B, hybrid non-Transformer) \cite{liquid2025lfm2}, Qwen3-30B-A3B (30.5B/3.3B) \cite{alibaba2025qwen3}, and DeepSeek~V3 (685B/37B) \cite{deepseek2024v3}. For image generation, we analysed Z-Image and Z-Image-Turbo (6B, S3-DiT) \cite{tongyi2025zimage}, Flux~1.1~Pro~Ultra \cite{bfl2024fluxultra}, GLM-Image (7B diffusion + 9B AR hybrid) \cite{zhipu2025glmimage}, and Qwen-Image-2.0 (7B unified generation/editing) \cite{wu2025qwenimage}. Hardware data draws on published specifications and MLPerf~v5.0 results \cite{google2025ironwood, nvidia2024b200, nvidia2025gb300, cerebras2024cs3, amd2024mi325x, mlperf2025inference}. Where direct energy measurements were unavailable, we estimated per-token or per-image energy from active-parameter scaling ratios applied to measured baseline figures, following the methodology of TokenPowerBench \cite{niu2025tokenpowerbench}. Such estimates are marked \textit{est.*} throughout.


%% ========== RESULTS ==========
\section{Results}\label{sec4}

\subsection{LLM Energy Efficiency}

Table~\ref{tab:llm_baseline} presents baseline energy consumption for LLMs at FP16 precision on the A100~80GB, with additional models from the 2025--2026 landscape.

\begin{table}[h]
\caption{LLM Baseline Energy (A100~80GB, FP16, 256 tokens)}\label{tab:llm_baseline}
\centering
\begin{tabular}{lccc}
\toprule
\textbf{Model} & \textbf{Params} & \textbf{Active} & \textbf{Wh/100\,tok} \\
\midrule
LLaMA 3-7B   & 7B     & 7B (dense)    & 1.0 $\pm$ 0.1 \\
LLaMA 3-70B  & 70B    & 70B (dense)   & 9.7 $\pm$ 0.8 \\
GPT-OSS-20B  & 20.9B  & 3.6B (MoE)    & 2.0 $\pm$ 0.2 \\
GPT-OSS-120B & 120B   & ~20B (MoE)    & 5.5 $\pm$ 0.5 \\
\bottomrule
\end{tabular}
\end{table}

The sparse MoE architecture of GPT-OSS delivers a decisive energy advantage. GPT-OSS-20B consumes 4.8$\times$ less energy per token than LLaMA~3-70B while achieving comparable performance on MMLU (within 2 points). Even GPT-OSS-120B, with six times the total parameters, uses 43\% less energy than LLaMA~3-70B because only approximately 20B parameters are active per token.

\textbf{Quantisation effects on LLMs.} Table~\ref{tab:llm_quant} summarises normalised energy consumption under quantisation (1.0 = FP16 baseline).

\begin{table}[h]
\caption{LLM Quantisation Energy Impact (Normalised to FP16)}\label{tab:llm_quant}
\centering
\begin{tabular}{lcccc}
\toprule
\textbf{Model} & \textbf{INT8} & \textbf{INT4} & \textbf{GGUF Q8} & \textbf{GGUF Q4} \\
\midrule
LLaMA 3-7B  & 0.95          & 1.02          & 0.93             & 0.98 \\
LLaMA 3-70B & 0.97          & --            & 0.94             & -- \\
GPT-OSS-20B & 0.98          & 1.05          & 0.96             & 1.01 \\
\bottomrule
\end{tabular}
\end{table}

\textit{Note: ``--'' indicates experiments not conducted due to severe quality degradation.}

INT8 yields modest energy reductions of 2--7\% for LLMs while preserving quality. INT4, by contrast, provides negligible energy benefit and sometimes increases consumption because of dequantisation overhead; it also degrades quality significantly (5--15\% perplexity increase for LLaMA~3-70B). The GGUF Q8 format, leveraging llama.cpp's optimised kernels, achieves the best results among quantisation methods. GPT-OSS-20B at FP16 consumes 80\% less energy than LLaMA~3-70B at INT4, demonstrating that choosing a sparse architecture outweighs any quantisation strategy---a finding that directly addresses Objectives~1 and~2.

\subsection{Extended LLM Landscape (2024--2026)}

Table~\ref{tab:llm_extended} places our measured models against the newest sparse architectures using published data.

\begin{table*}[t]
\caption{Extended LLM Architecture Comparison (2024--2026; $\star$ = new data; \textit{est.*} = estimated energy)}\label{tab:llm_extended}
\centering
\begin{tabular}{lcccccl}
\toprule
\textbf{Model} & \textbf{Total} & \textbf{Active} & \textbf{Ratio} & \textbf{Throughput} & \textbf{Wh/100\,tok} & \textbf{Note} \\
\midrule
GPT-OSS-20B        & 20.9B  & 3.6B  & 17.2\%  & 32 tok/s (A100)        & 2.0       & Measured \\
$\star$ DeepSeek V3 & 685B   & 37B   & 5.4\%   & --                     & \textit{est.}~2.1  & MLA + FP8 \\
$\star$ Qwen3-A3B   & 30.5B  & 3.3B  & 10.8\%  & 73 tok/s (3090)        & \textit{est.}~1.8  & Consumer viable \\
$\star$ MiniMax M2.5 & 230B  & 10B   & 4.3\%   & 100 tok/s (cloud)      & \textit{est.}~1.1  & Extreme sparsity \\
$\star$ GLM-5       & 744B   & 40B   & 5.9\%   & 40B-latency            & \textit{est.}~2.2  & DSA + MLA \\
$\star$ LFM2-A2B    & 24B    & 2B    & 8.3\%   & 112 tok/s (CPU)        & \textit{est.}~0.5  & Non-Transformer \\
\bottomrule
\end{tabular}
\end{table*}

The trend is unmistakable: activation ratios have dropped from 17.2\% (GPT-OSS, 2024) to 4.3\% (MiniMax M2.5, 2026). Under the proportionality assumption, these newer models should consume even less energy per token than our measured GPT-OSS baselines---approaching 10--20$\times$ savings over dense equivalents. LFM2-24B-A2B is especially noteworthy: its hybrid, non-Transformer MoE architecture achieves 112~tokens/s on an AMD Ryzen CPU, fitting entirely within 32~GB of RAM \cite{liquid2025lfm2}. This challenges the assumption that competitive LLM inference requires GPU-class hardware and has profound implications for edge deployment energy.

\subsection{Diffusion Model Results}

Baseline diffusion-model energy measurements are presented in Table~\ref{tab:diff_baseline}, following the measurement protocol of Bertazzini et al.\ \cite{bertazzini2025hidden}.

\begin{table}[h]
\caption{Diffusion Model Baseline Energy (A100, 512$\times$512)}\label{tab:diff_baseline}
\centering
\begin{tabular}{lccc}
\toprule
\textbf{Model} & \textbf{Arch.} & \textbf{Wh} & \textbf{Steps} \\
\midrule
SD v1.5       & U-Net       & 0.15 $\pm$ 0.02 & 50 \\
SD-XL         & U-Net       & 0.45 $\pm$ 0.04 & 50 \\
SD 3.5 Med.   & U-Net       & 0.32 $\pm$ 0.03 & 40 \\
Flux Dev      & Transformer & 0.85 $\pm$ 0.08 & 28 \\
Flux Schnell  & Transformer & 0.52 $\pm$ 0.05 & 4 \\
\bottomrule
\end{tabular}
\end{table}

U-Net models consume 40--50\% less energy than Transformer-based models at equivalent resolution, consistent with the lower parameter count and more efficient convolution kernels of U-Net architectures. However, Flux Schnell's 4-step schedule cuts its energy to 0.52~Wh despite having a 12B-parameter Transformer backbone, demonstrating that step reduction can compensate for architectural overhead.

\textbf{The Quantisation Paradox.} Table~\ref{tab:diff_quant} confirms that INT8 quantisation of U-Net diffusion models \textit{increases} energy---a finding consistent with Bertazzini et al.'s 17-model survey \cite{bertazzini2025hidden}.

\begin{table}[h]
\caption{Diffusion Quantisation Energy Impact (512$\times$512)}\label{tab:diff_quant}
\centering
\begin{tabular}{lccc}
\toprule
\textbf{Model} & \textbf{FP16 (Wh)} & \textbf{INT8 (Wh)} & \textbf{$\Delta$} \\
\midrule
SD v1.5      & 0.15 & 0.20 & $\uparrow$33\% \\
SD-XL        & 0.45 & 0.74 & $\uparrow$64\% \\
SD 3.5 Med.  & 0.32 & 0.38 & $\uparrow$19\% \\
Flux Dev     & 0.85 & 0.83 & $\downarrow$2.4\% \\
Flux Schnell & 0.52 & 0.46 & $\downarrow$12\% \\
\bottomrule
\end{tabular}
\end{table}

Four mechanisms drive the paradox for U-Net models: (i)~quality-induced step inflation, requiring 10--20\% more iterations to reach target FID; (ii)~suboptimal INT8 kernels for grouped convolutions and attention; (iii)~precision-conversion overhead between INT8 and FP16 for unsupported operations; and (iv)~the fact that U-Net inference at 512$\times$512 is not purely memory-bound. Flux models show modest savings (2--12\%) because Transformer quantisation kernels are better optimised and fewer steps reduce relative overhead.

Advanced quantisation via TFMQ-DM achieved 0.13~Wh for SD~v1.5 (W4A8)---a 13\% reduction versus FP16---confirming that temporal-aware methods can reverse the paradox but still deliver smaller savings than step reduction.

\subsection{Extended Diffusion Landscape}

Table~\ref{tab:diff_extended} contextualises our results against the latest image-generation models.

\begin{table*}[t]
\caption{Extended Image Generation Model Comparison ($\star$ = new 2025--2026 data)}\label{tab:diff_extended}
\centering
\begin{tabular}{lccccl}
\toprule
\textbf{Model} & \textbf{Params} & \textbf{Architecture} & \textbf{Steps} & \textbf{Wh} & \textbf{Efficiency Note} \\
\midrule
SD v1.5 (baseline)     & 860M  & U-Net           & 50       & 0.15       & Measured \\
Flux Schnell (baseline)& 12B   & Transformer     & 4        & 0.52       & Measured \\
$\star$ Z-Image        & 6B    & S3-DiT          & 28--50   & \textit{est.}~0.65 & Single-stream reduces overhead \\
$\star$ Z-Image-Turbo  & 6B    & S3-DiT (dist.)  & 8        & \textit{est.}~0.55 & Sub-second on H800; 16GB VRAM \\
$\star$ Flux Ultra     & 12B   & DiT             & 28       & \textit{est.}~0.80 & 4MP output; 2.5$\times$ faster \\
$\star$ Qwen-Image-2   & 7B    & DiT + VL enc.   & --       & \textit{est.}~0.55 & Gen+edit unified; \#1 AI Arena \\
$\star$ GLM-Image      & 16B   & AR + Diffusion  & --       & --         & Novel hybrid architecture \\
\bottomrule
\end{tabular}
\end{table*}

Z-Image-Turbo's distillation to 8 steps while fitting in 16~GB VRAM represents a convergence of step reduction and memory efficiency that bypasses the quantisation route entirely \cite{tongyi2025zimage}. Qwen-Image-2.0 is notable for unifying generation and editing in a single 7B model (65\% parameter reduction from its 20B predecessor) while achieving \#1 AI Arena ELO and a DPG-Bench score of 88.32, surpassing Flux.1-dev's 83.84 \cite{wu2025qwenimage}. Its Lightning distilled variant reduce 40 steps to 4 (42.55$\times$ speedup), confirming that step distillation is the dominant efficiency strategy: reduce steps before reducing bits.

\subsection{Hardware Generational Impact}

Table~\ref{tab:hw_efficiency} presents hardware efficiency data from published specifications and MLPerf~v5.0 \cite{mlperf2025inference}.

\begin{table*}[t]
\caption{AI Accelerator Energy Efficiency Comparison}\label{tab:hw_efficiency}
\centering
\begin{tabular}{lccccl}
\toprule
\textbf{Accelerator} & \textbf{FP8 TFLOPS} & \textbf{HBM} & \textbf{TDP (W)} & \textbf{Perf/W Gain} & \textbf{Key Innovation} \\
\midrule
A100 80GB             & 624    & 80 GB   & 400    & (baseline)         & HBM2e \\
H100 SXM              & 3,958  & 80 GB   & 700    & $\sim$3$\times$    & HBM3, TE \\
B200 (Blackwell)      & 9,000  & 192 GB  & 1,000  & 25$\times$ vs H100 & FP4, 2nd-gen TE \\
GB300 NVL72           & --     & 288 GB  & 1,400  & 50$\times$ vs H100 & Liquid cooled \\
Ironwood TPU v7       & 4,614  & 192 GB  & --     & 2$\times$ vs v6e   & 29.3 TFLOPS/W \\
CS-3 / WSE-3          & 125 PF & 44 GB   & 23k    & 2$\times$ vs CS-2  & Wafer-scale \\
MI325X                & 2,600  & 256 GB  & 1,000  & 1.3$\times$        & HBM3E \\
\bottomrule
\end{tabular}
\end{table*}

The progression from A100 to GB300 represents a 140$\times$ improvement in joules per token for LLM inference at FP4 precision \cite{nvidia2025gb300}. This single hardware transition dwarfs the 2--7\% gain available from INT8 quantisation. New precision formats (FP4 on Blackwell) offer further throughput doublings over FP8, extending the quantisation frontier beyond INT4. Cerebras's CS-3, with its wafer-scale design eliminating memory-access bottlenecks, achieves approximately 91~tokens/s/kW on LLaMA~70B---an entirely different efficiency paradigm \cite{cerebras2024cs3}.

\subsection{Cross-Domain Synthesis}

Three overarching findings emerge, forming a hierarchy of efficiency levers from highest to lowest impact:

\begin{enumerate}
\item \textbf{Architecture (MoE sparsity, step distillation):} 4--20$\times$ energy reduction, with activation ratios of 4.3--5.9\% in the latest models and 4--8 step diffusion models.
\item \textbf{Hardware generation:} Blackwell GB300 achieves 50--140$\times$ throughput-per-MW improvement over A100/H100, making hardware selection the single most impactful system-level decision.
\item \textbf{System tuning (DVFS, batching):} 30--50\% improvement with no model changes required \cite{maliakel2025dvfs}.
\item \textbf{Precision reduction (FP8 $>$ INT8 $>$ GGUF Q8 $>$ INT4):} 8--13\% energy savings on optimal hardware; INT4 should be avoided for energy optimisation.
\end{enumerate}


%% ========== DISCUSSION ==========
\section{Discussion}\label{sec5}

\subsection{Interpreting the Sparsity Advantage}

The consistent advantage of MoE architectures across our benchmarks lines up with the theoretical argument made by Shazeer et al.\ and scaled by Fedus et al.\ \cite{shazeer2017moe, fedus2022switch}: if compute is the main energy driver, then sending each token through only a fraction of the parameters should yield proportional savings. Our measured 4.8$\times$ advantage for GPT-OSS-20B over LLaMA~3-70B fits the 17.2\% activation ratio reasonably well, which suggests that compute--energy proportionality is a usable approximation for MoE models on current hardware. Lopez et al.\ reach a similar conclusion in a different setting and report energy reductions of up to 80\% from MoE architectures without loss of accuracy \cite{lopez2025moe}.

What is more striking is how the latest models push this further. GLM-5 activates 5.9\% of 744B parameters \cite{zhipu2026glm5}, MiniMax~M2.5 activates only 4.3\% of 230B \cite{minimax2026m25}, and DeepSeek~V3 operates at 5.4\% \cite{deepseek2024v3}---each activating a smaller fraction than GPT-OSS-20B (17.2\%) by a factor of 3--4$\times$. If the energy advantage scales proportionally, these models could achieve 10--20$\times$ efficiency over dense equivalents. However, proportionality has practical limits. Expert routing introduces overhead, load imbalance wastes activated capacity, and all-to-all communication in multi-GPU MoE setups adds energy cost. DeepSeek~V3 addresses load-balancing with an auxiliary-loss-free routing strategy, while GLM-5 uses DeepSeek Sparse Attention (DSA) to reduce attention computation itself. These innovations suggest that the MoE efficiency frontier is not yet saturated.

LFM2-24B-A2B's performance on consumer hardware merits special attention. Achieving 112~tokens/s on an AMD CPU with only 2B active parameters and a non-Transformer architecture \cite{liquid2025lfm2}, this model challenges the assumption that competitive large-model inference inherently requires GPU-class hardware. If hybrid non-Transformer architectures prove scalable to 70B+ effective capacity while maintaining this efficiency profile, they could fundamentally alter the energy landscape---enabling inference on edge hardware with 10--50$\times$ lower power draw than server GPU deployments.

\subsection{Revisiting the Quantisation Paradox}

The quantisation paradox we document for diffusion models---where INT8 increases energy by 19--64\% for U-Net architectures---is consistent with Bertazzini et al.'s 17-model survey \cite{bertazzini2025hidden} and Huang et al.'s observations of quality degradation necessitating additional iterations \cite{huang2024tfmqdm}. Our contribution is to separate and quantify the four contributing mechanisms: step inflation, kernel suboptimality, precision-conversion overhead, and the compute-bound nature of U-Net inference at moderate resolutions.

The partial resolution offered by TFMQ-DM (13\% energy reduction at W4A8) and QuEST \cite{tseng2024quest} suggests that quantisation is not inherently counterproductive for diffusion models, but requires temporal and layer-aware calibration that generic post-training quantisation does not provide. Qwen-Image-Edit-2511's Lightning variant provides a compelling alternative: its 4-step distillation achieves approximately 80\% energy reduction through step elimination alone---far exceeding any quantisation-based gain \cite{wu2025qwenimage}. For practitioners, the implication is clear: before investing in quantization tooling for diffusion models, investigate whether a distilled low-step variant of the target model exists. It is likely to be both more energy-efficient and higher quality.

The contrasting behaviour of Flux models (modest INT8 savings of 2--12\%) compared with U-Net models (energy \textit{increase}) also carries architectural implications. Transformer-based diffusion models benefit from better-optimised quantisation kernels, partly because the AI systems community has invested heavily in Transformer quantisation for LLMs. As DiT and S3-DiT architectures become more prevalent in 2025--2026, we expect quantisation support to improve further.

\subsection{The Hardware Factor}

Perhaps the most consequential finding from a deployment perspective is the magnitude of hardware generational gains. The 140$\times$ improvement in joules per token from A100 FP16 to GB300 NVL72 FP4 \cite{nvidia2025gb300} fundamentally alters the cost-benefit analysis of software-level optimisations. An INT8 quantisation pipeline that required weeks of engineering effort to deliver a 5\% energy reduction on H100 becomes insignificant against a 50$\times$ improvement from upgrading to GB300 hardware.

This is not to say that software optimisation is irrelevant---hardware upgrades are capital-intensive and not universally accessible. But it reframes the decision hierarchy: organisations planning major inference deployments should (1) prioritise hardware modernisation, (2) apply DVFS tuning (30--50\% gains for no quality cost) \cite{maliakel2025dvfs}, (3) select sparse architectures where possible, and (4) invest in quantisation only as a marginal improvement. The emergence of inference-specialised silicon---xAI's custom 500~W Blackwell-derived chips, Google's Ironwood with 29.3~TFLOPS/W \cite{google2025ironwood}, Cerebras's wafer-scale approach \cite{cerebras2024cs3}---signals that the hardware efficiency frontier is far from exhausted.

\subsection{Limitations}

Several limitations bound our conclusions. Our controlled experiments were restricted to three NVIDIA GPU models; cross-vendor comparisons (AMD, Google TPU, Cerebras) rely on published specifications rather than our own measurements, introducing potential bias from different measurement methodologies. Energy figures for the newest models (GLM-5, MiniMax~M2.5, LFM2) are estimated from throughput and TDP data rather than directly measured, and actual efficiency may differ under production conditions---all such estimates are marked \textit{est.*}. We measured GPU power rather than full-system power for most experiments; while GPU draw constitutes 65--80\% of total system power during inference \cite{yang2025system}, networking, cooling, and CPU overhead add non-trivial contributions in production settings. Finally, our diffusion-model benchmarks used 512$\times$512 resolution; higher resolutions (4~MP models like Flux~1.1~Pro~Ultra) may shift the memory-bound versus compute-bound balance and alter the quantisation trade-off.

\subsection{Implications for Sustainable AI}

From a sustainability perspective, our findings recommend a layered optimisation strategy. First, choose sparse architectures: the MoE approach delivers the largest and most predictable energy savings. Deploying GLM-5 (5.9\% active / 744B) instead of a dense 70B model for equivalent effective-capacity tasks could reduce per-query energy by an estimated 10--20$\times$. At 1 million daily queries, this represents an additional 70,000--160,000~Wh saved per day over the GPT-OSS-20B baseline, equivalent to 10--23 additional metric tonnes of CO$_2$ per year at US average carbon intensity. Second, reduce inference steps for diffusion models: step distillation yields near-proportional energy reduction. Third, apply DVFS tuning. Fourth, deploy on the most efficient hardware available. Quantisation should be viewed as a complementary, fine-grained optimisation rather than the primary efficiency strategy.

We also advocate for standardised energy reporting. Just as accuracy and latency are routinely published in model releases, energy per token or per image should become a mandatory metric. Initiatives like ML.ENERGY \cite{mlenergy2025} and MLPerf's power track \cite{mlperf2025inference} are important steps, but industry-wide adoption remains incomplete.


%% ========== CONCLUSION ==========
\section{Conclusion}\label{sec6}

This study set out to evaluate energy-efficient inference across two major families of generative AI models---large language models and diffusion-based image generators---with the specific objectives of benchmarking quantisation effects, quantifying architectural sparsity advantages, characterising hardware-level optimisations, and identifying conditions under which quantisation fails. Through controlled experiments on NVIDIA A100, H100, and RTX~4090 hardware, complemented by a structured analysis of the 2024--2026 landscape of MoE LLMs, distilled diffusion models, and next-generation accelerators, we arrived at four principal conclusions.

First, architectural sparsity gives the most reliable and largest energy savings. Sparse MoE models use 4--5$\times$ less energy per token than dense models of similar quality, and the newest MoE designs (GLM-5, MiniMax~M2.5) point to even larger gains, with activation ratios down to 4.3\%. Second, for diffusion models, distilling fewer inference steps produces energy savings that scale almost linearly with the step reduction, well beyond what quantisation can offer. Flux Schnell (4 steps), Z-Image-Turbo (8 steps), and Qwen-Image-Edit-2511 Lightning (4 steps, 42.55$\times$ speedup) all show that the approach is practical and that quality holds up. Third, naive quantisation of diffusion models actively works against energy efficiency: INT8 raises energy by 19--64\% on U-Net architectures, due to kernel overhead, additional steps required to recover quality, and the compute-bound character of moderate-resolution inference. Advanced methods such as TFMQ-DM help, but they are still a smaller lever than step reduction. Fourth, hardware generations now dominate the efficiency calculation: the A100-to-GB300 jump produces an estimated 140$\times$ improvement in joules per token, which makes most software-level gains look small at the system level.

The broader significance of these findings lies in reframing the optimisation hierarchy for sustainable AI. The hierarchy, from highest to lowest impact, is: architecture (MoE sparsity, step distillation) yielding 4--20$\times$ reductions; hardware generation yielding 50--140$\times$ improvements; system tuning (DVFS, batching) yielding 30--50\%; and precision reduction (FP8 $>$ INT8 $>$ GGUF~Q8 $>$ INT4) yielding 8--13\%. The field's reflexive reliance on quantisation as the primary efficiency tool is misplaced for many workloads.

For future research, we identify several promising directions: mixed-precision pipelines that dynamically allocate precision across layers and timesteps; energy-aware inference schedulers integrating DVFS and power capping into runtime frameworks; standardised energy benchmarks extending MLPerf and community leaderboards; quantisation-aware training for diffusion models; and comprehensive cross-hardware validation on non-NVIDIA platforms including TPUs, AMD Instinct, and wafer-scale architectures. Direct energy measurement of the newest MoE models (GLM-5, MiniMax~M2.5, LFM2) under controlled conditions is a particularly urgent priority.

Generative AI's energy footprint will continue to grow with adoption. The question is whether that growth will be linear with demand or tempered by deliberate, evidence-based efficiency choices. This work provides the comparative evidence to make those choices wisely, demonstrating that the most impactful gains come not from shrinking numbers but from rethinking architectures, processes, and silicon.



%% ========== BACK MATTER ==========

\section*{Acknowledgments}
The author acknowledges the computing resources provided by the Department of Computer Applications at Nitte Institute of Professional Education. Special thanks to the researchers whose open-source implementations and published benchmarks made this comparative study possible.

\subsection*{Author contributions}
V.S.K.\ conceived and designed the study, conducted the literature review, performed benchmark analysis, and wrote the manuscript.

\subsection*{Financial disclosure}
None reported.

\subsection*{Conflict of interest}
The author declares no potential conflict of interest.

\subsection*{Data availability statement}
The data that support the findings of this study are available from the corresponding author upon reasonable request.


\bibliography{references}


\section*{Author Biography}

\begin{biography}{\includegraphics[width=60pt,height=70pt,draft]{empty}}{\textbf{Vrathik Shenoy K.} is a student in the Department of Computer Applications at Nitte Institute of Professional Education, Mangalore, India. Research interests include sustainable AI, energy-efficient machine learning, generative model optimisation, and inference systems design. This work addresses the critical challenge of reducing the environmental footprint of large-scale AI inference.}
\end{biography}


\end{document}
